% 第四章 研究方法
\chapter{研究方法}\label{ch:method}

\section{尾數歸零與三種捨入方式}\label{sec:rounding}

對一個有 $p$ 個尾數位元的運算元（FP64 為 52、BF16 為 7、FP8 E4M3 為 3），保留最高的 $k$ 個位元，其餘 $p-k$ 個設為 0。本研究比較三種去除低位的方式：

\begin{description}
  \item[截斷：] 直接以位元遮罩清除低位，即往 0 的方向捨去。
  \item[最近捨入：] 先加上保留的最低位權重的一半（即被捨去部分的最高位），再以遮罩清除。攔截工具、GEMM 微基準、Quantum ESPRESSO、DF-MP2 與高斯過程的實作在平手時往離零較遠的方向捨入；HPL、LoRA 的權重載入與語言模型推論的實作則是平手取偶。
  \item[隨機捨入：] 先加上一個小於被捨去位元權重的隨機整數，再以遮罩清除。
\end{description}

每個元素只需一次整數 AND 運算（最近捨入再多一次加法，隨機捨入另需產生亂數），可以融入產生該運算元的核心；本報告的 HPL 另外執行一個捨入核心、攔截工具另外做一次複製，這些額外成本都計入量到的時間。本報告所稱「截斷到 $k$ 位」泛指保留 $k$ 位；除非特別註明，應用程式（Quantum ESPRESSO、PySCF、DF-MP2、高斯過程、ExaGeoStat 與第三方程式）與 HPL 使用最近捨入，GEMM 微基準、分界點掃描與 HPL-MxP 使用往零截斷。HPL 的位數依 rocHPL 的慣例為含隱藏位的有效位數：36 位相當於保留 35 個尾數位元。各實驗的位數定義、去除低位的方式與樣本數整理於表~\ref{tab:settings}。

\begin{table}[htbp]
\centering
\caption{各實驗的位數定義、去除低位的方式與樣本數}\label{tab:settings}
\small\renewcommand{\arraystretch}{1.3}
\begin{tabularx}{\textwidth}{@{}>{\raggedright\arraybackslash}p{3.6cm}>{\raggedright\arraybackslash}p{2.9cm}>{\raggedright\arraybackslash}p{3.1cm}L@{}}
\toprule
實驗 & $k$ 的意義 & 去除低位的方式 & 樣本 \\
\midrule
GEMM 微基準、機制對照、分界點掃描 & 尾數位元 & 往零截斷 & 微基準 3 節點 $\times$ 2 對（95\% 區間）；機制 1000~W 每個組合 2–4 次；分界點掃描每格兩輪 \\ \addlinespace[2pt]
HPL（8 節點） & 有效位數，\newline 含隱藏位\newline（36 = 35 個\newline 尾數位元） & 最近捨入\newline（平手取偶） & 兩組節點，每個上限 2 + 4 對（第二組先暖機），95\% 區間 \\ \addlinespace[2pt]
HPL-MxP（8 節點） & BF16 尾數位元 & 往零截斷 & 兩組節點，每個上限 2 + 4 對（第二組先暖機），95\% 區間 \\ \addlinespace[2pt]
Quantum ESPRESSO、PySCF、\newline DF-MP2、高斯過程、\newline ExaGeoStat、MFEM、\newline SW4、LAMMPS & 尾數位元 & 最近捨入\newline（平手離零） & QE 先暖機一次再交錯 6 對；DF-MP2 與高斯過程每格 6 輪；以上附 95\% 區間；其餘多為每格兩輪（位數掃描與攔截版的效能對照各一次） \\ \addlinespace[2pt]
LoRA 訓練 & FP8 尾數位元 & 往零截斷\newline（另測最近捨入與隨機捨入） & 兩個 seed、四組配對 \\ \addlinespace[2pt]
語言模型推論 & BF16/FP8\newline 尾數位元 & 最近捨入\newline（平手取偶） & 六個模型，各一次 \\
\bottomrule
\end{tabularx}
\end{table}

\section{四種實作層次}\label{sec:impl}

\begin{enumerate}
  \item \textbf{在基準程式中歸零。}矩陣乘法微基準在產生輸入時歸零；HPL 在每次 FP64 GEMM 之前以一個捨入核心處理兩個輸入（執行時以環境變數指定位數）；HPL-MxP 在編譯時指定位數，每次寫回 BF16 時歸零。
  \item \textbf{攔截函式庫。}對不修改原始碼的科學軟體，以 \texttt{LD\_PRELOAD} 攔下 cuBLAS 的所有 FP64 與複數 FP64 矩陣乘法入口，把 A、B 複製並捨入後再呼叫原函式（呼叫端的資料不變）。同一個函式庫也能統計運算元的指數分佈、改用 cuBLAS 的 FP64 模擬，或改以 FP32/3xTF32 計算，用於公平比較。
  \item \textbf{程式內改寫。}在 Quantum ESPRESSO 的 Davidson 對角化中，對波函數區塊原地歸零，不需要複製；第三方非矩陣乘法程式（MFEM、SW4、LAMMPS）在設定階段對儲存的係數陣列歸零一次；GPU4PySCF 的 DF-MP2 在三指標張量建好後原地歸零一次；高斯過程核心則在每個面板進行尾端更新之前原地歸零。
  \item \textbf{依收斂門檻自動調整。}Davidson 每一圈會收到一個收斂門檻 ethr，在整個 SCF 過程中從約 $10^{-2}$ 收緊到約 $10^{-13}$。我們保留
  \begin{equation}
    k = \min\bigl(52,\ \max\bigl(16,\ \lceil -\log_2 \mathrm{ethr} \rceil + m\bigr)\bigr)
    \label{eq:adapt}
  \end{equation}
  個位元，其中 $m$ 為餘裕，使前期使用較少位元、最後幾圈才用足。這是一條經驗規則：ethr 是特徵值的絕對容許誤差（單位 Ry），與尾數的相對誤差不同，$m$ 用來吸收兩者的差距；此規則不需要對應用建模，只使用程式本來就會計算的容許誤差。
\end{enumerate}

LoRA 訓練則在載入後對凍結的 FP8 基底權重歸零一次；推論實驗在線性層中模擬各種格式，其中 FP8 使用真實的 FP8 tensor core 矩陣乘法。

\section{實驗平台}\label{sec:platform}

實驗平台為 14 個節點的叢集，每個節點有一顆 NVIDIA GH200 480GB superchip（72 核 Grace CPU、Hopper GPU 搭配 96~GB HBM3，最高 SM 時脈 1980~MHz），以 200~Gb/s InfiniBand 連接，由 Slurm 管理。模組功耗上限可設為 200–1000~W（預設 1000~W）。規格表列出的可設範圍是 450–1000~W~\cite{NVIDIAGH200}；本報告低於 450~W 的設定超出規格表的範圍，但驅動程式接受，實測功耗也貼著上限。軟體為 NVHPC 26.5（CUDA 13.2、cuBLAS 13、NCCL、HPC-X Open MPI）。多節點實驗（HPL、HPL-MxP、LoRA）使用 8 個節點。完整的軟體版本列於附錄 \ref{app:env}。

第二個平台是兩個各有 8 張 AMD Instinct MI350X 的節點（ROCm 7.2.3），功耗上限為每張 GPU 的封裝功耗上限（1–1000~W，預設 1000~W），不含 CPU。它的平台細節、量測方式與統計單位與 GH200 不同，在第 \ref{sec:mi-setup} 節說明。

\section{量測與統計}\label{sec:measure}

\begin{description}
  \item[功耗：] 每 250–500~ms 以 nvidia-smi 取樣模組功耗、SM 時脈、溫度與降頻原因。耗電量以驅動程式的 1 秒平均模組功耗對時間積分（對忽高忽低的負載，單一瞬時讀數會有 $-37\%$ 到 $+41\%$ 的偏差）；表中的平均功耗是量測視窗內瞬時取樣的平均，穩態下與 1 秒平均相差約 1\% 以內。
  \item[公平比較：] 每次量測都等 GPU 降到 45\,°C 以下才開始；在同一功耗上限內，以交錯順序（A、B、B、A）在同一組節點上執行各變體，以消除溫度與先後順序的影響。
  \item[信賴區間：] 同一節點（多節點基準為同一組節點）、同一上限、同一輪的「歸零 ÷ 原生 $-$ 1」為一個成對樣本，以 $t$ 分布取 95\% 區間，藉此消除節點間 13–17\% 的時脈差異。矩陣乘法微基準為 3 個節點各 2 對；Quantum ESPRESSO 為先暖機後的 6 對；DF-MP2 與高斯過程為兩個工作共 6 輪；HPL 與 HPL-MxP 為兩組節點共 6 對（第一組 2 對、第二組先暖機後 4 對）。區間把同一工作或同一組節點內的各對當成獨立樣本，而它們並非完全獨立，所以區間可能偏窄。每個工作的第一次執行常偏慢（暖機），而交錯順序的第一個多半是原生，會使第一對的收益偏大，所以 Quantum ESPRESSO 的量測，以及 DF-MP2、高斯過程、HPL 與 HPL-MxP 的第二個工作，都先執行一次不計入的暖機（第 \ref{sec:sci} 節）。
  \item[MI350X：] 以硬體的能量累加器計算能量，信賴區間以真正獨立的單位（量測視窗、整節點的一次執行或 8 張 GPU 的一輪）計算，並以溫度迴歸與受控的熱歷史處理風扇造成的偏誤（第 \ref{sec:mi-setup}、\ref{sec:mi-temp} 節）。
  \item[資料檢查：] 每個量測視窗都記錄實際生效的功耗上限，與預定值不符者一律剔除（叢集為共用資源）；硬體功耗煞車與溫度降頻以降頻計數器判讀（在 HPL 的兩次掃描 3315、3582 與 HPL-MxP 的診斷工作 1675 中開啟），或以取樣時記錄的降頻原因判讀。
\end{description}

\section{等效功耗 \texorpdfstring{$\Delta W$}{ΔW}}\label{sec:dw}

在功耗上限 $P$ 下，將尾數歸零版本的效能記為 $p'(P)$，原生版本的效能記為 $p(P)$。沿著原生版本實測的「效能對上限」曲線線性內插，找出原生版本達到 $p'(P)$ 所需的上限 $P_{\mathrm{native}}$，並定義
\begin{equation}
  \Delta W(P) = P_{\mathrm{native}} - P, \qquad p(P_{\mathrm{native}}) = p'(P) .
  \label{eq:dw}
\end{equation}
其意義為：未修改的程式要追上尾數歸零版本，每個節點需要多給的功耗上限。這比百分比更直接地對應到機房的電力預算。上限之間以直線內插；由於效能對上限的曲線向下彎，兩個上限相隔較遠時（例如 800–1000~W），$\Delta W$ 會高估。

\section{功耗上限下的效能模型}\label{sec:powermodel}

等效功耗需要在每個上限下量測每個截斷版本。以下的模型只用原生程式的效能曲線、平台的閒置功耗與一兩次量測，就能預測截斷版本在每個上限下的效能。設 GPU 在上限 $P$ 下停在目標時脈 $f$，電壓 $V(f)$ 由目標時脈決定，但切換發生在有效時脈 $sf$，其中 $s\le1$ 是隨資料而變的時脈降低（MI350X 上觀察到，第 \ref{sec:mi-mech} 節；沒有時 $s=1$）。以 $c$ 表示程式的等效切換電容、$P_0$ 表示不隨切換變化的功耗，則
\begin{equation}
  P = P_0 + c\,s\,V(f)^2 f = P_0 + c\,s\,g(f),
  \label{eq:pmodel}
\end{equation}
效能則與 $sf$ 成正比。截斷改變 $c$ 與 $s$，不改變 $g$。對原生（$n$）與截斷（$r$）版本定義
\begin{equation}
  R = \frac{c_r s_r}{c_n s_n}, \qquad \sigma = \frac{s_r}{s_n} ;
  \label{eq:rsigma}
\end{equation}
在相同的鎖定目標時脈下，$R$ 就是兩者動態功耗之比，$\sigma$ 是兩者有效時脈之比，都可以直接量測。在上限 $P$ 下，截斷版本停在 $c_r s_r g(f_r) = P - P_0$ 的目標時脈 $f_r$；原生版本要達到同樣的 $f_r$，需要上限 $P'$，滿足 $c_n s_n g(f_r) = P' - P_0$。兩式相除得
\begin{equation}
  P' = P_0 + \frac{P-P_0}{R}, \qquad p_r(P) = \sigma\, p_n(P'), \qquad
  \Delta W = P' - P = (P-P_0)\Bigl(\frac{1}{R}-1\Bigr),
  \label{eq:pmodel-pred}
\end{equation}
其中 $p_r(P)=\sigma\,p_n(P')$ 是因為兩者都在目標時脈 $f_r$ 運作，而效能與 $s$ 成正比；這裡的 $\Delta W$ 與式~\eqref{eq:dw} 的定義相同。模型假設 $s$ 在所用的時脈範圍內不隨 $f$ 改變、程式維持計算受限；預測值不超過程式受時脈限制的峰值。

模型需要原生的效能曲線 $p_n(P)$ 與 $P_0$。我們把 $P_0$ 設為實測的閒置功耗，與驗證資料無關：GH200 模組為 174~W（三個節點 169–179~W），MI350X 每張 GPU 為 262~W。截斷版本以 1300~MHz 的鎖頻量測取得 $R$ 與 $\sigma$（MI350X，不使用上限掃描的任何一點），或以最低上限的速度推得 $R$（GH200 無法鎖頻，$\sigma=1$，該點不列入評估）。驗證結果見第 \ref{sec:mi-model} 節。式~\eqref{eq:pmodel-pred} 也說明了等效功耗 $\Delta W\propto P-P_0$ 隨上限增加，而相對加速並不一定如此。

\section{誤差模型}\label{sec:errmodel}

兩個輸入都保留 $k$ 位時，矩陣乘法結果的相對誤差（Frobenius 範數）約為
\begin{equation}
  \frac{\lVert \widehat{C} - C \rVert_F}{\lVert C \rVert_F} \approx \alpha \cdot 2^{-k},
  \label{eq:errmodel}
\end{equation}
與內積長度無關（累加仍為 FP64）。在均勻隨機資料上（$k \ge 4$），最近捨入的 $\alpha \approx 0.27$、隨機捨入 0.38、截斷 0.70，截斷的誤差約為最近捨入的 2.6 倍。截斷的誤差幾乎全是系統性的縮小，對會累積偏差的流程有害，這是科學軟體應採用最近捨入的原因（第 \ref{sec:sci} 節）。
